\section{Proof of Theorem \ref{Theorem:Main1}}\label{Section:Proof_Main1}
We show that $\hat{f}(x)={}_r\varphi_s(\bm{a};\bm{b};q,x)$ is $[\lambda;p]$-summable for any $\lambda\in\mathbb{C}\setminus\big[(-1)^k;q\big]$. Let $g(\xi)$ be the $p$-Borel transform of $\hat{f}(x)$
\begin{gather*}
g(\xi)=(\hat{\mathcal{B}}_{p;1}\hat{f})(\xi)=\sum_{n\ge0}\frac{(a_1,a_2,\ldots,a_r;q)_n}{(b_1,\ldots,b_s;q)_n(q;q)_n}\big\{(-1)^nq^{\frac{n(n-1)}{2}}\big\}^{1+s-r}p^{\frac{n(n-1)}{2}}\xi^n.
\end{gather*}
Then, $g(\xi)$ is convergent in $|\xi|<1$ and again can be expressed by basic hypergeometric series:
\begin{gather*}
g(\xi)=\sum_{n\ge0}\frac{(a_1,a_2,\ldots,a_r;q)_n}{(b_1,\ldots,b_s;q)_n(q;q)_n}\big((-1)^{-k}\xi\big)^n={}_r\varphi_{r-1}\left(
\begin{matrix}
a_1,a_2,\ldots,a_r\\
b_1,b_2,\ldots,b_s,0,\ldots,0
\end{matrix}
;q;(-1)^{-k}\xi
\right).
\end{gather*}
To have the analytic continuation of $g(\xi)$, we use the following proposition.
\begin{Proposition}[Slater \cite{Slater1952} or Gaspar--Rahman {\cite[Section~4.5]{Gasper&Rahman2004}}] The analytic continuation of the basic hypergeometric series ${}_r\varphi_{r-1}(a_1,a_2,\ldots,a_r;b_1,\ldots,b_{r-1};q,\xi)$ is given by
\begin{gather}
{}_r\varphi_{r-1}\left(
\begin{matrix}a_1,\ldots,a_r\\
b_1,\ldots,b_{r-1}
\end{matrix};q,\xi\right)=\frac{(a_2,\ldots,a_r,b_1/a_1,\ldots,b_{r-1}/a_1;q)_\infty}{(b_1,\ldots,b_{r-1},a_2/a_1,\ldots,a_r/a_1;q)_\infty}\frac{\theta_q(-a_1\xi)}{\theta_q(-\xi)}\nonumber\\
\hphantom{{}_r\varphi_{r-1}\left(
\begin{matrix}a_1,\ldots,a_r\\
b_1,\ldots,b_{r-1}
\end{matrix};q,\xi\right)=}{} \times{}_r\varphi_{r-1}\left(
\begin{matrix}a_1,a_1q/b_1,\ldots,a_1q/b_{r-1}\\
a_1q/a_2,\ldots,a_1q/a_r
\end{matrix};q,\frac{qb_1\cdots b_{s}}{a_1\cdots a_r\xi }\right)\nonumber\\
\hphantom{{}_r\varphi_{r-1}\left(
\begin{matrix}a_1,\ldots,a_r\\
b_1,\ldots,b_{r-1}
\end{matrix};q,\xi\right)=}{} +\idem(a_1;a_2,\ldots,a_r).\label{r,r-1AC}
\end{gather}
\end{Proposition}
Letting $b_{s+1},\ldots,b_{r-1}\to 0$ in \eqref{r,r-1AC}, it is seen that $g(\xi)$ is continued analytically to $\xi\in\mathbb{C}^*\setminus\big[(-1)^{k};q\big]$ as follows:
\begin{gather}
g(\xi)=\frac{(a_2,\ldots,a_r,b_1/a_1,\ldots,b_{s}/a_1;q)_\infty}{(b_1\ldots,b_s,a_2/a_1,\ldots a_r/a_1;q)_\infty}\frac{\theta_q((-1)^{k-1}a_1\xi)}{\theta_q((-1)^{k-1}\xi)}\nonumber\\
\hphantom{g(\xi)=}{} \times{}_{s+1}\varphi_{r-1}\left(\tilde{\bm{a}};\tilde{\bm{b}};q,(-1)^k\frac{q^{k+1}b_1\cdots b_s}{a_1^{1-k}a_2\cdots a_r\xi}\right)\nonumber\\
\hphantom{g(\xi)=}{} +\idem(a_1;a_2,\ldots,a_r),\label{ac_borelim}
\end{gather}
where $\tilde{\bm{a}}=(a_1,a_1q/b_1,\ldots,a_1q/b_s)\in\mathbb{C}^{s+1}$ and $\tilde{\bm{b}}=(a_1q/a_2,\ldots,a_1q/a_r)\in\mathbb{C}^{r-1}$. We note that the basic hypergeometric series ${}_{s+1}\varphi_{r-1}$ in~\eqref{ac_borelim} is holomorphic on $\mathbb{C}^*$. In order to obtain \eqref{ac_borelim}, we use
\begin{gather*}
\lim_{b\to0}(aq/b;q)_nb^n=\lim_{b\to0}(b-aq)\big(b-aq^2\big)\cdots\big(b-aq^n\big)=(-a)^nq^{\frac{n(n+1)}{2}}.
\end{gather*}

Now, we shall show that $\big(\hat{\mathcal{B}}_{p;1}\hat{f}\big)(\xi)=g(\xi)\in\mathbb{H}_{p;1}^{[\lambda;p]}$ for any $\lambda\in\mathbb{C}^*\setminus\big[(-1)^k;q\big]$. The following lemma holds.
\begin{Lemma}\label{Lemma:Main1}
The function $g(\xi)$ has a $p$-exponential growth at infinity in
\begin{gather*}
\Omega_\delta:=\mathbb{C}^*\setminus\bigcup_{m\in\mathbb{Z}}\big\{\xi\in\mathbb{C}^* ;\big\vert\xi-(-1)^{k}q^m\vert<\delta\big\vert q^m\vert\big\}
\end{gather*}
for any $\delta>0$.
\end{Lemma}
A proof of this lemma will be given in Section~\ref{Section:Proof_Lemma:Main1}. For any $\lambda\in\mathbb{C}^*\setminus\big[(-1)^k;q\big]$, there exists positive constants $\delta$ and $\varepsilon$ such that
\begin{gather*}
\bigcup_{m\in\mathbb{Z}}\big\{\xi\in\mathbb{C}^* ;\,\vert\xi-\lambda p^m\vert<\varepsilon\vert p^m\lambda\vert\big\}\subset\Omega_\delta.
\end{gather*}
Therefore we obtain $g(\xi)\in\mathbb{H}_{p;1}^{[\lambda;p]}$ for any $\lambda\in\mathbb{C}^*\setminus\big[(-1)^k;q\big]$. In the other words, $\hat{f}(x)={}_r\varphi_s(\bm{a};\bm{b};q,x)$ is $[\lambda;p]$-summable for any $\lambda\in\mathbb{C}^*\setminus\big[(-1)^k;q\big]$.

We consider the $p$-Laplace transform of $g(\xi)$. For any $\lambda\in\mathbb{C}^*\setminus\big[(-1)^k;q\big]$, let ${}_rf_s(\bm{a};\bm{b};\lambda;q,x)$ be the $p$-Laplace transform of $g(\xi)$
\begin{gather*}
{}_rf_s(\bm{a};\bm{b};\lambda;q,x)=\big(\mathcal{L}_{p;1}^{[\lambda;p]}g\big)(x)=\sum_{m\in\mathbb{Z}}\frac{g\big(\lambda p^m\big)}{\theta_p\big(\frac{\lambda p^m}{x}\big)}.
\end{gather*}
Since
\begin{gather*}
\frac{\theta_q\big((-1)^{k-1}a_1\lambda p^m\big)}{\theta_q\big((-1)^{k-1}\lambda p^m\big)}=\frac{\theta_q\big((-1)^{k-1}a_1\lambda q^{km}\big)}{\theta_q\big((-1)^{k-1}\lambda q^{km}\big)}=\frac{1}{a_1^{km}}\frac{\theta_q\big((-1)^{k-1}a_1\lambda\big)}{\theta_q\big((-1)^{k-1}\lambda\big)},
\\
{}_{s+1}\varphi_{r-1}\left(\tilde{\bm{a}};\tilde{\bm{b}};q,(-1)^k\frac{q^{k+1}b_1\cdots b_s}{a_1^{1-k}a_2\cdots a_r\lambda p^m}\right)=\sum_{n\ge0}\frac{(\tilde{\bm{a}};q)_n}{(\tilde{\bm{b}},q;q)_n}p^{\frac{n(n-1)}{2}}\left(\frac{pqb_1\cdots b_s}{a_1^{1-k}a_2\cdots a_r\lambda p^m}\right)^n,
\end{gather*}
and
\begin{gather*}
\theta_p\left(\frac{\lambda p^m}{x}\right)=\left(\frac{\lambda}{x}\right)^{-m}p^{-\frac{m(m-1)}{2}}\theta_p\left(\frac{\lambda}{x}\right),
\end{gather*}
we have
\begin{gather*}
{}_rf_s(\bm{a};\bm{b};\lambda;q,x)= \frac{(a_2,\ldots,a_r,b_1/a_1,\ldots,b_{s}/a_1;q)_\infty}{(b_1\ldots,b_s,a_2/a_1,\ldots, a_r/a_1;q)_\infty}\frac{\theta_q\big((-1)^{k-1}a_1\lambda\big)}{\theta_q\big((-1)^{k-1}\lambda\big)}\frac{1}{\theta_p(\lambda/x)}\\
\hphantom{{}_rf_s(\bm{a};\bm{b};\lambda;q,x)=}{} \times\sum_{n\ge0}\frac{(\tilde{\bm{a}};q)_n}{(\tilde{\bm{b}},q;q)_n}\left(\frac{qb_1\cdots b_s}{a_1^{1-k}a_2\cdots a_r\lambda}\right)^n \sum_{m\in\mathbb{Z}}\left(\frac{\lambda}{a_1^kx}\right)^mp^{\frac{n(n-1)}{2}+n-mn+\frac{m(m-1)}{2}}\\
\hphantom{{}_rf_s(\bm{a};\bm{b};\lambda;q,x)=}{} +\idem(a_1;a_2,\ldots,a_r).
\end{gather*}
Here we remark that
\begin{gather*}
\sum_{m\in\mathbb{Z}}\left(\frac{\lambda}{a_1^kx}\right)^mp^{\frac{n(n-1)}{2}+n-mn+\frac{m(m-1)}{2}} =\left(\frac{\lambda}{a_1^kx}\right)^n\sum_{m\in\mathbb{Z}}p^{\frac{(m-n)(m-n-1)}{2}}\left(\frac{\lambda}{a_1^kx}\right)^{m-n}\\
\hphantom{\sum_{m\in\mathbb{Z}}\left(\frac{\lambda}{a_1^kx}\right)^mp^{\frac{n(n-1)}{2}+n-mn+\frac{m(m-1)}{2}} }{}
=\left(\frac{\lambda}{a_1^kx}\right)^n\theta_p\left(\frac{\lambda}{a_1^kx}\right).
\end{gather*}
Therefore we obtain
\begin{gather*}
{}_rf_s(\bm{a};\bm{b};\lambda;q,x)= \frac{(a_2,\ldots,a_r,b_1/a_1,\ldots,b_{s}/a_1;q)_\infty}{(b_1\ldots,b_s,a_2/a_1,\ldots a_r/a_1;q)_\infty}\frac{\theta_q\big((-1)^{k-1}a_1\lambda\big)}{\theta_q\big((-1)^{k-1}\lambda\big)}\frac{\theta_p\big(\lambda/a_1^kx\big)}{\theta_p(\lambda/x)}\\
\hphantom{{}_rf_s(\bm{a};\bm{b};\lambda;q,x)=}{} \times\sum_{n\ge0}\frac{(\tilde{\bm{a}};q)_n}{(\tilde{\bm{b}},q;q)_n}\left(\frac{qb_1\cdots b_s}{a_1a_2\cdots a_rx}\right)^n\\
\hphantom{{}_rf_s(\bm{a};\bm{b};\lambda;q,x)=}{} +\idem(a_1;a_2,\ldots,a_r).
\end{gather*}
From
\begin{gather*}
\frac{\theta_p\big(\lambda/a_1^kx\big)}{\theta_p(\lambda/x)}=\frac{\theta_p\big(pa_1^kx/\lambda\big)}{\theta_p(px/\lambda)}
\end{gather*}
(is seen from equality \eqref{theta_inverse}) and
\begin{gather*}
\sum_{n\ge0}\frac{(\tilde{\bm{a}};q)_n}{(\tilde{\bm{b}},q;q)_n}\left(\frac{qb_1\cdots b_s}{a_1a_2\cdots a_rx}\right)^n={}_r\varphi_{r-1}\left(
\begin{matrix}\tilde{\bm{a}},\bm{0}_k\\
\tilde{\bm{b}}
\end{matrix};q,\frac{qb_1\cdots b_s}{a_1a_2\cdots a_rx}\right),
\end{gather*}
we finish the proof.
\section{Proof of Lemma \ref{Lemma:Main1}}\label{Section:Proof_Lemma:Main1}
In this section, we give a proof of Lemma \ref{Lemma:Main1}, that is, we show
\begin{gather*}
\vert g(\xi)\vert\le L \theta_{\vert p\vert}(M\vert \xi\vert),\qquad \xi\in\Omega_\delta,
\end{gather*}
where $L$ and $M$ are positive constants.

As easily seen, the following inequality holds.
\begin{gather*}
\theta_{|p|}(A|\xi|)+\theta_{|p|}(B|\xi|)\le \theta_{|p|}(C|\xi|),
\end{gather*}
where $A, B>0$ and $C=A+B$. Hence, showing that the first term of \eqref{ac_borelim} has $p$-exponential growth at infinity in $\Omega_\delta$ is sufficient to prove Lemma \ref{Lemma:Main1}. We write the first term of \eqref{ac_borelim} as~$g_1(\xi)$, i.e.,
\begin{gather*}
g_1(\xi):=\frac{(a_2,\ldots,a_r,b_1/a_1,\ldots,b_{s}/a_1;q)_\infty}{(b_1\ldots,b_s,a_2/a_1,\ldots a_r/a_1;q)_\infty}\frac{\theta_q\big((-1)^{k-1}a_1\xi\big)}{\theta_q\big((-1)^{k-1}\xi\big)}\\
\hphantom{g_1(\xi):=}{}\times {}_{s+1}\varphi_{r-1}\left(\tilde{\bm{a}};\tilde{\bm{b}};q;(-1)^k\frac{q^{k+1}b_1\cdots b_s}{a_1^{1-k}a_2\cdots a_r\xi}\right),
\end{gather*}
where $\tilde{\bm{a}}=(a_1,a_1q/b_1,\ldots,a_1q/b_s)\in\mathbb{C}^{s+1}$ and $\tilde{\bm{b}}=(a_1q/a_2,\ldots,a_1q/a_r)\in\mathbb{C}^{r-1}$.

We show that basic hypergeometric series ${}_{s+1}\varphi_{r-1}$ has $p$-exponential growth. It holds that
\begin{gather*}
|(\alpha;q)_n| =|1-\alpha||1-\alpha q|\cdots\big|1-\alpha q^{n-1}\big|\le(1+|\alpha|)^n,\\
|(a_iq/a_j;q)_n| =|1-a_iq/a_j|\big|1-a_iq^2/a_j\big|\cdots\big|1-a_iq^n/a_j\big|\ge \left(\inf_{k\in\mathbb{N}^*}\big\{\big|1-a_iq^k/a_j\big|\big\}\right)^n>0
\end{gather*}
for any $\alpha\in\mathbb{C}$ and $1\le i,j\le r$. Therefore we obtain
\begin{align*}
\left|{}_{s+1}\varphi_{r-1}\left(\tilde{\bm{a}};\tilde{\bm{b}};q;(-1)^k\frac{q^{k+1}b_1\cdots b_s}{a_1^{1-k}a_2\cdots a_r\xi}\right)\right|&\le\sum_{n\ge0}\frac{|(\tilde{\bm{a}};q)_n|}{|(\tilde{\bm{b}},q;q)_n|}|p|^{\frac{n(n+1)}{2}}\left|\frac{qb_1\cdots b_s}{a_1^{1-k}a_2\cdots a_r\xi}\right|^n\\
&\le\sum_{n\ge0}|p|^{\frac{n(n+1)}{2}}\left(\frac{1}{M|\xi|}\right)^n\\
&\le\sum_{n\in\mathbb{Z}}|p|^{\frac{n(n+1)}{2}}\left(\frac{1}{M|\xi|}\right)^n=\theta_{|p|}(M|\xi|),
\end{align*}
where $M$ is a positive constant. Hence we have
\begin{gather}\label{g1_estimate_1}
|g_1(\xi)|\le C\left|\frac{\theta_q\big((-1)^{k-1}a_1\xi\big)}{\theta_q\big((-1)^{k-1}\xi\big)}\right|\theta_{|p|}(M|\xi|).
\end{gather}
Here $C$ is a positive constant.

Now, for any $\xi \in\Omega_\delta$, there exist an integer $n$ and
\begin{gather*}
\xi_0\in\Omega_\delta \cap\{\xi\in\mathbb{C}^* ;\,|p|\le|\xi|\le1\}
\end{gather*}
such that $\xi=p^n\xi_0$.
Substituting $\xi=p^n\xi_0$ into \eqref{g1_estimate_1}, we have
\begin{gather*}
|g_1(\xi)|\le C\left|\frac{\theta_q((-1)^{k-1}a_1\xi_0)}{\theta_q((-1)^{k-1}\xi_0)}\right|\big(M|a_1|^k\big)^{-n}|p|^{-\frac{n(n-1)}{2}}\theta_{|p|}(M|\xi_0|).
\end{gather*}
Since
\begin{gather*}
\big(M|a_1|^k\big)^{-n}|p|^{-\frac{n(n-1)}{2}}=\frac{\theta_{|p|}\big(M|a_1|^k|p|^n|\xi_0|\big)}{\theta_{|p|}\big(M|a_1|^k|\xi_0|\big)}=\frac{\theta_{|p|}\big(M|a_1|^k|\xi|\big)}{\theta_{|p|}\big(M|a_1|^k|\xi_0|\big)}
\end{gather*}
holds, we obtain
\begin{gather*}
|g_1(\xi)|\le C\left|\frac{\theta_q\big((-1)^{k-1}a_1\xi_0\big)}{\theta_q\big((-1)^{k-1}\xi_0\big)}\right|\frac{\theta_{|p|}(M|\xi_0|)}{\theta_{|p|}\big(M|a_1|^k|\xi_0|\big)}\theta_{|p|}\big(M|a_1|^k|\xi|\big).
\end{gather*}
Now, the function of $\xi_0$
\begin{gather*}
\left|\frac{\theta_q\big((-1)^{k-1}a_1\xi_0\big)}{\theta_q\big((-1)^{k-1}\xi_0\big)}\right|\frac{\theta_{|p|}(M|\xi_0|)}{\theta_{|p|}\big(M|a_1|^k|\xi_0|\big)}
\end{gather*}
is holomorphic on the compact set $\Omega_\delta \cap\{\xi\in\mathbb{C}^* ;\,|p|\le|\xi|\le1\}$. Hence there exists a positive constant $L$ such that
\begin{gather*}
\max_{\Omega_\delta \cap\{\xi\in\mathbb{C}^* ;\,|p|\le|\xi|\le1\}}\left|\frac{\theta_q\big((-1)^{k-1}a_1\xi_0\big)}{\theta_q\big((-1)^{k-1}\xi_0\big)}\right|\frac{\theta_{|p|}(M|\xi_0|)}{\theta_{|p|}\big(M|a_1|^k|\xi_0|\big)}\le L.
\end{gather*}

